\chapter{Manual de código y repositorio}
\label{anexo:manual_codigo}

Este anexo actúa únicamente como punto de referencia para el acceso al código
fuente del proyecto. Harmony Hub se entrega como repositorio versionado, con la
aplicación móvil, el backend, la documentación académica y los artefactos de
validación integrados en una única estructura de trabajo.

La versión pública y consultable del proyecto se encuentra en:

\begin{center}
\url{https://github.com/i22remua/THE-HARMONY-HUB}
\end{center}

Desde ese repositorio pueden revisarse la organización del cliente Flutter, la
implementación del backend FastAPI, la integración con Spotify, el catálogo
musical estructurado, el componente de aprendizaje y la propia memoria del
TFG.

\section*{Puertas de calidad del modelo supervisado}

Para evitar que el componente de aprendizaje actúe con evidencia insuficiente,
Harmony Hub incorpora una política explícita de puertas de calidad antes de
permitir que el modelo reordene candidatos de sesión. Esta explicación resume
lo desarrollado en el diseño de la solución y facilita localizarlo en el código
y en los artefactos de auditoría.

En primer lugar, se aplica una \textbf{puerta de datos}. El modelo solo puede
considerarse operativo cuando existe un mínimo de ejemplos totales y una
cobertura mínima por clase. En términos prácticos, esto significa que no basta
con acumular muchas sesiones, sino que deben existir suficientes casos útiles y
no útiles para que el clasificador vea ambos comportamientos. En la
implementación actual se exige \texttt{total\_examples >= 40} y al menos
\texttt{8} ejemplos por clase.

En segundo lugar, cuando la distribución lo permite, se utiliza
\textbf{validación cruzada estratificada}. Esta técnica divide el conjunto de
ejemplos en varios pliegues conservando proporciones similares de clases en
cada partición. Su objetivo es reducir el riesgo de considerar válido un modelo
que solo funciona bien en una única separación afortunada de los datos.

Sobre esa validación se calculan varias métricas. La
\textbf{balanced accuracy} mide si el modelo diferencia de forma razonable las
sesiones útiles y las no útiles sin sesgarse hacia la clase mayoritaria. El
\textbf{ROC AUC} mide si el clasificador sabe asignar más probabilidad a los
casos positivos que a los negativos antes incluso de fijar un umbral binario.
La métrica \textbf{F1-score} resume el equilibrio entre precisión y recobrado
de la clase positiva.

Para sintetizar estas señales se utiliza un \textbf{índice compuesto de
calidad}, calculado como:

\begin{center}
\texttt{quality score = (0.45 * balanced accuracy + 0.35 * ROC AUC + 0.20 * F1) * 100}
\end{center}

En la configuración actual, el modelo solo supera la puerta de calidad si
alcanza simultáneamente:

\begin{itemize}
    \item \texttt{balanced accuracy >= 0.45},
    \item \texttt{ROC AUC >= 0.38},
    \item \texttt{quality score >= 49}.
\end{itemize}

Además, incluso cuando el modelo ha sido entrenado y auditado correctamente, su
influencia sobre una recomendación concreta queda condicionada por una
\textbf{puerta de confianza operacional}. Esta última exige que la probabilidad
del mejor candidato supere un umbral mínimo antes de modificar el ranking
heurístico. En el estado actual del proyecto, ese valor operativo aparece
registrado como \texttt{min\_selected\_mode\_probability = 0.14}.

Estas comprobaciones pueden localizarse directamente en los siguientes puntos
del repositorio:

\begin{itemize}
    \item \texttt{backend\_fastapi/app/services/session\_mode\_ml\_service.py}: umbrales operativos y activación del modelo en tiempo de recomendación.
    \item \texttt{backend\_fastapi/app/ml/train\_ranking\_model.py}: entrenamiento, validación cruzada, cálculo de métricas y decisión de paso de puertas.
    \item \texttt{backend\_fastapi/app/ml/models/session\_mode\_model\_audit.json}: evidencia persistida del estado de auditoría y de los umbrales vigentes.
    \item \texttt{backend\_fastapi/app/ml/models/session\_mode\_model\_card.json}: resumen funcional del modelo, su propósito y sus parámetros de operación.
\end{itemize}

Esta política de control tiene una función importante dentro del proyecto: el
aprendizaje no sustituye a la base heurística, sino que solo interviene cuando
dispone de evidencia mínima, calidad suficiente y confianza operacional en el
contexto evaluado.
